\section{Exercício Prático: App com Persistência}

\begin{frame}
  \frametitle{Exercício: Persistência de Dados}
  \textbf{Objetivo}: Criar um Dockerfile para uma aplicação Node.js, gerar a imagem e
  garantir que arquivos enviados pelo usuário não sumam ao reiniciar o container.
\end{frame}

\begin{frame}[fragile]
  \frametitle{Dockerfile da Aplicação}
\begin{verbatim}
FROM node:22-alpine
WORKDIR /app
COPY package.json ./
RUN npm install --production
COPY server.js ./
RUN mkdir -p /app/uploads
EXPOSE 3000
CMD ["node", "server.js"]
\end{verbatim}

  \vspace{0.3em}
  \texttt{docker build -t minha-app .}
\end{frame}

\begin{frame}[fragile]
  \frametitle{Sem Volume: Dados Perdidos}
\begin{verbatim}
docker run -d -p 3000:3000 \
  --name app-sem-volume minha-app
# ... enviar arquivos ...
docker rm -f app-sem-volume
# Arquivos enviados são perdidos!
\end{verbatim}
\end{frame}

\begin{frame}[fragile]
  \frametitle{Com Volume: Dados Persistem}
\begin{verbatim}
docker volume create uploads-data

docker run -d -p 3000:3000 --name app-com-volume \
  -v uploads-data:/app/uploads minha-app

# ... enviar arquivos ...
docker rm -f app-com-volume

# Novo container, mesmo volume
docker run -d -p 3000:3000 --name app-restaurado \
  -v uploads-data:/app/uploads minha-app

# Arquivos ainda estão lá!
\end{verbatim}
\end{frame}

\begin{frame}
  \frametitle{Conclusão}
  \begin{block}{Regra}
    Identificar quais diretórios contêm dados que devem sobreviver ao ciclo
    de vida do container (uploads, bancos, logs) e montar volumes nesses diretórios.
  \end{block}

  \vspace{0.5em}
  Mesmo que o container seja parado, removido ou recriado, os dados permanecem intactos.
\end{frame}
